\frfilename{mt263.tex}
\versiondate{4.4.13}
\copyrightdate{2000}

\def\diam{\mathop{\text{diam}}}
\def\dist{\mathop{\text{dist}}}

\def\chaptername{Perubahan variabel dalam integral}
\def\sectionname{Transformasi terdiferensialkan dalam $\BbbR^r$}

\newsection{263}

\def\headlinesectionname{Transformasi terdiferensialkan dalam
$\eightBbb R^r$}

Bagian ini dikhususkan untuk pembuktian satu teorema utama
(263D) mengenai transformasi terdiferensialkan antara himpunan-himpunan bagian
$\BbbR^r$.   Hasil ini akan digeneralisasi dalam \S265,
dan pembaca yang cukup mengenal pokok bahasan ini, atau cukup berani,
mungkin ingin membaca \S264 sebelum mempelajari bagian ini dan \S265
secara bersamaan.   Saya mengakhiri bagian ini dengan beberapa akibat sederhana dan suatu
perluasan hasil utama yang dapat dibuat dalam kasus berdimensi satu (263J).

Di seluruh bagian ini, sebagaimana dalam bagian lain bab ini, $\mu$ akan
menyatakan ukuran Lebesgue pada $\BbbR^r$.

\leader{263A}{Transformasi \dvrocolon{linear}}\cmmnt{ Saya mulai dengan
kasus khusus
operator linear, yang bukan saja menjadi dasar pembuktian
263D, melainkan juga salah satu penerapannya yang paling penting,
dan bahkan sudah memadai untuk banyak hasil yang sangat mencolok.

\medskip

\noindent}{\bf Teorema} Misalkan $T$ suatu matriks real $r\times r$;  pandang
$T$ sebagai
operator linear dari $\BbbR^r$ ke dirinya sendiri.   Misalkan $J=|\det T|$ adalah
modulus determinannya.   Maka

\Centerline{$\mu T[E]=J\mu E$}

\noindent untuk setiap himpunan terukur $E\subseteq\BbbR^r$.   Jika
$T$ merupakan suatu permutasi (yakni, jika $J\ne 0$), maka

\Centerline{$\mu F=J\mu(T^{-1}[F])$}

\noindent untuk setiap himpunan terukur $F\subseteq\BbbR^r$, dan

\Centerline{$\int_F g\,d\mu=J\int_{T^{-1}[F]} gT\,d\mu$}

\noindent untuk setiap fungsi terintegralkan $g$ dan himpunan terukur $F$.

\proof{{\bf (a)} Langkah pertama ialah menunjukkan bahwa $T[I]$
terukur untuk setiap interval setengah-terbuka $I\subseteq\BbbR^r$.   \Prf\
Setiap interval setengah-terbuka tak kosong
$I=\coint{a,b}$ adalah gabungan terhitung dari interval-interval tertutup
$I_n =[a,b-2^{-n}\tbf{1}]$, dan setiap $I_n$ kompak (2A2F),
sehingga $T[I_n]$ kompak (2A2Eb), maka tertutup (2A2Ec),
maka terukur (115G), dan $T[I]=\bigcup_{n\in\Bbb N}T[I_n]$
terukur.\ \Qed

\medskip

{\bf (b)} Tetapkan $J^*=\mu T[\,\coint{\tbf{0},\tbf{1}}\,]$, dengan
$\tbf{0}=(0,\dots,0)$ dan $\tbf{1}=(1,\ldots,1)$;  karena
$T[\,\coint{\tbf{0},\tbf{1}}\,]$ terbatas, $J^*<\infty$.   (Pada akhirnya saya akan
menunjukkan bahwa $J^*=J$.)   Lebih mudah untuk terlebih dahulu menangani kasus
$T$ yang singular.   Ingat bahwa $T$, jika dipandang sebagai transformasi linear
dari $\BbbR^r$ ke dirinya sendiri, bersifat bijektif atau mempunyai citra berupa
suatu subruang linear sejati.   Dalam kasus yang kedua, ambil sebarang $e\in\Bbb
R^r\setminus T[\BbbR^r]$;  maka himpunan-himpunan

\Centerline{$T[\,\coint{\tbf{0},\tbf{1}}\,]+\gamma e$,}

\noindent ketika $\gamma$ menjelajahi $[0,1]$, saling lepas dan semuanya
berukuran sama, yaitu $J^*$, karena $\mu$ invarian terhadap translasi (134A);
selain itu, gabungannya terbatas, sehingga mempunyai ukuran luar berhingga.   Karena
terdapat tak berhingga banyak $\gamma$ semacam itu, ukuran bersama $J^*$ haruslah
nol.   Sekarang perhatikan bahwa

\Centerline{$T[\BbbR^r]=\bigcup_{z\in \Bbb
Z^r}T[\,\coint{\tbf{0},\tbf{1}}\,]+Tz$,}

\noindent dan

\Centerline{$\mu(T[\,\coint{\tbf{0},\tbf{1}}\,]+Tz)=J^*=0$}

\noindent untuk setiap $z\in\Bbb Z^r$, sedangkan $\Bbb Z^r$ terhitung, sehingga
$\mu T[\Bbb R^r]=0$.   Pada saat yang sama, karena
$T$ singular, determinannya nol, dan $J=0$.
Dengan demikian

\Centerline{$\mu T[E]=0=J\mu E$}

\noindent untuk setiap himpunan terukur $E\subseteq\BbbR^r$, dan pembuktian selesai.

\medskip

{\bf (c)} Mulai sekarang, dengan demikian, kita andaikan bahwa $T$
nonsingular.
Perhatikan bahwa operator ini dan inversnya kontinu, sehingga $T$ merupakan
homeomorfisme, dan $T[G]$ terbuka jika dan hanya jika $G$ terbuka.

Jika $a\in\BbbR^r$ dan $k\in\Bbb N$, maka

\Centerline{$\mu T[\,\coint{a,a+2^{-k}\tbf{1}}\,]=2^{-kr}J^*$.}

\noindent   \Prf\ Tetapkan
$J^*_k=\mu T[\,\coint{\tbf{0},2^{-k}\tbf{1}}\,]$.   Sekarang
$T[\,\coint{a,a+2^{-k}\tbf{1}}\,]
=T[\,\coint{\tbf{0},2^{-k}\tbf{1}}\,]+Ta$;  karena $\mu$
invarian terhadap translasi, ukurannya juga $J^*_k$.   Selanjutnya,
$\coint{\tbf{0},\tbf{1}}$ dapat dinyatakan sebagai gabungan saling lepas dari $2^{kr}$
himpunan berbentuk $\coint{a,a+2^{-k}\tbf{1}}$;  akibatnya,
$T[\,\coint{\tbf{0},\tbf{1}}\,]$ dapat dinyatakan sebagai gabungan saling lepas dari
$2^{kr}$ himpunan berbentuk $T[\,\coint{a,a+2^{-k}\tbf{1}}\,]$, dan


\Centerline{$J^*=\mu T[\,\coint{\tbf{0},\tbf{1}}\,]=2^{kr}J^*_k$,}

\noindent yakni,
$J^*_k=2^{-kr}J^*$, sebagaimana diklaim.\ \Qed

\medskip

{\bf (d)} Akibatnya $\mu T[G]=J^*\mu G$ untuk setiap himpunan terbuka
$G\subseteq\BbbR^r$.   \Prf\ Untuk setiap $k\in\Bbb N$, tetapkan

\Centerline{$Q_k
=\{z:z\in\Bbb Z^r,\,\coint{2^{-k}z,2^{-k}z+2^{-k}\tbf{1}}\subseteq G\}$,}

\Centerline{$G_k=\bigcup_{z\in Q_k}
\coint{2^{-k}z,2^{-k}z+2^{-k}\tbf{1}}$.}

\noindent Maka $G_k$ adalah gabungan saling lepas dari $\#(Q_k)$ himpunan berbentuk
$\coint{2^{-k}z,2^{-k}z+2^{-k}\tbf{1}}$, sehingga $\mu G_k=2^{-kr}\#(Q_k)$;
juga, $T[G_k]$ adalah gabungan saling lepas dari $\#(Q_k)$ himpunan berbentuk
$T[\,\coint{2^{-k}z,2^{-k}z+2^{-k}\tbf{1}}\,]$, sehingga berukuran
$2^{-kr}J^*\#(Q_k)=J^*\mu G_k$, dengan menggunakan (c).

Selanjutnya perhatikan bahwa $\sequence{k}{G_k}$ adalah barisan tak menurun dengan
gabungan $G$, sehingga

\Centerline{$\mu T[G]=\lim_{k\to\infty}\mu T[G_k]
=\lim_{k\to\infty}J^*\mu G_k
=J^*\mu G$.   \Qed}

\medskip

{\bf (e)} Maka $\mu^*T[A]=J^*\mu^* A$ untuk setiap
$A\subseteq\Bbb R^r$.   \Prf\ Diberikan $A\subseteq\BbbR^r$ dan
$\epsilon>0$, terdapat
himpunan-himpunan terbuka $G$, $H$ sedemikian sehingga $G\supseteq A$, $H\supseteq T[A]$,
$\mu G\le\mu^* A+\epsilon$ dan $\mu H\le\mu^*T[A]+\epsilon$ (134Fa).
Tetapkan $G_1=G\cap T^{-1}[H]$;  maka $G_1$ terbuka karena $T^{-1}[H]$
terbuka.   Sekarang $\mu T[G_1]=J^*\mu G_1$, sehingga

$$\eqalign{\mu^* T[A]
&\le\mu T[G_1]
=J^*\mu G_1
\le J^*\mu^* A+J^*\epsilon\cr
&\le J^*\mu G_1+J^*\epsilon
=\mu T[G_1]+J^*\epsilon
\le\mu H+J^*\epsilon\cr
&\le\mu^*T[A]+\epsilon+J^*\epsilon.\cr}$$

\noindent Karena $\epsilon$ sebarang, $\mu^*T[A]=J^*\mu^* A$.\ \Qed

\medskip

{\bf (f)} Akibatnya $\mu T[E]$ ada dan sama dengan $J^*\mu E$
untuk setiap himpunan terukur $E\subseteq\BbbR^r$.   \Prf\ Misalkan
$E\subseteq\BbbR^r$ terukur, dan ambil sebarang
$A\subseteq\BbbR^r$.   Tetapkan $A'=T^{-1}[A]$.   Maka

$$\eqalign{\mu^*(A\cap T[E])+\mu^*(A\setminus T[E])
&=\mu^*(T[A'\cap E])+\mu^*(T[A'\setminus E])\cr
&=J^*(\mu^*(A'\cap E)+\mu^*(A'\setminus E))\cr
&=J^*\mu^* A'
=\mu^*T[A']
=\mu^* A.\cr}$$

\noindent Karena $A$ sebarang, $T[E]$ terukur, dan sekarang


\Centerline{$\mu T[E]=\mu^*T[E]=J^*\mu^* E=J^*\mu E$.  \Qed}

\medskip

{\bf (g)} Akhirnya kita siap menghitung $J^*$.
Ingat bahwa matriks $T$ harus dapat dinyatakan sebagai $PDQ$, dengan $P$ dan
$Q$ matriks ortogonal dan $D$ diagonal, dengan entri-entri diagonal
nonnegatif (2A6C).  Maka harus berlaku

\Centerline{$T[\,\coint{\tbf{0},\tbf{1}}\,]
=P[D[Q[\,\coint{\tbf{0},\tbf{1}}\,]]]$,}

\noindent sehingga, dengan menggunakan (f),

\Centerline{$J^*=J^*_PJ^*_DJ^*_Q$,}

\noindent dengan $J^*_P=\mu P[\,\coint{\tbf{0},\tbf{1}}\,]$, dan seterusnya.
Sekarang
kita mendapati bahwa $J^*_P=J^*_Q=1$.   \Prf\ Misalkan $B=B(\tbf{0},1)$ adalah bola
satuan dalam $\BbbR^r$.    Karena $B$ tertutup, ia terukur;  karena
terbatas, $\mu B<\infty$;  dan karena $B$ memuat interval setengah-terbuka tak kosong
$\coint{\tbf{0},r^{-1/2}\tbf{1}}$, $\mu B>0$.   Sekarang
$P[B]=Q[B]=B$, karena $P$ dan $Q$ adalah matriks ortogonal;
jadi kita mempunyai

\Centerline{$\mu B=\mu P[B]=J^*_P\mu B$,}

\noindent dan $J^*_P$ haruslah $1$;  demikian pula, $J^*_Q=1$.\ \Qed

\medskip

{\bf (h)} Jadi kita hanya perlu menghitung $J^*_D$.   Andaikan
koefisien-koefisien $D$ adalah $\delta_1,\ldots,\delta_r\ge 0$, sehingga
$Dx=(\delta_1\xi_1,\ldots,\delta_r\xi_r)\penalty-100=d\times x$.   Sejak
awal (c) kita telah mengandaikan bahwa $T$ nonsingular, sehingga
tidak ada $\delta_i$ yang bernilai $0$.   Dengan demikian

\Centerline{$D[\,\coint{\tbf{0},\tbf{1}}\,]
=\coint{\tbf{0},d}$,}

\noindent dan

\Centerline{$J^*_D=\mu\coint{\tbf{0},d}=\prod_{i=1}^r\delta_i
=\det D$.}

\noindent Sekarang karena $P$ dan $Q$ ortogonal, keduanya mempunyai determinan
$\pm 1$, sehingga $\det T=\pm\det D$ dan $J^*=\pm\det T$;  karena $J^*$ tentu
nonnegatif, $J^*=|\det T|=J$.

\medskip

{\bf (i)} Dengan demikian $\mu T[E]=J\mu E$ untuk setiap himpunan terukur Lebesgue
$E\subseteq\BbbR^r$.   Karena $T$ nonsingular, kita dapat menggunakan
argumen di atas untuk menunjukkan
bahwa $T^{-1}[F]$ terukur untuk setiap $F$ terukur, dan

\Centerline{$\mu F=\mu T[T^{-1}[F]]=J\mu T^{-1}[F]
=\int J\times\chi(T^{-1}[F])\,d\mu$,}

\noindent dengan mengidentifikasi $J$ dengan fungsi konstan bernilai $J$.
Berdasarkan 235A,

\Centerline{$\int_F g\,d\mu=\int_{T^{-1}[F]}JgT\,d\mu
=J\int_{T^{-1}[F]}gT\,d\mu$}

\noindent untuk setiap fungsi terintegralkan $g$ dan himpunan terukur $F$.
}%end of proof of 263A

\cmmnt{
\leader{263B}{Catatan} Mungkin saya seharusnya telah memperingatkan Anda bahwa saya
akan menggunakan hasil-hasil \S235.   Namun, jika hasil-hasil itu masih segar
dalam ingatan Anda, rumus-rumus dalam pernyataan teorema tentu telah
mengingatkannya; dan jika tidak, barangkali lebih baik kembali kepada hasil-hasil itu
sekarang daripada sebelum membaca teorema, karena hasil-hasil itu hanya digunakan dalam
kalimat terakhir pembuktian.

Saya telah menguraikan argumen di atas dengan santai,
bahkan mungkin dapat dikatakan sangat perlahan.   Intinya ialah bahwa meskipun
invariansi ukuran Lebesgue terhadap translasi, dan perilakunya di bawah
pembesaran sederhana pada satu koordinat, kurang lebih sudah tertanam dalam
definisinya, perilakunya di bawah rotasi umum tidak demikian, karena suatu
rotasi membawa interval setengah-terbuka menjadi balok miring.   Tentu saja
perhitungan ukuran objek semacam itu sebenarnya tidak berkaitan khusus
dengan teori Lebesgue, dan jelas bahwa sebagian besar argumen ini akan berlaku
sama baiknya bagi setiap pengertian volume berdimensi $r$ yang wajar secara geometris.

Sekarang kita sampai pada hasil utama bab ini.  Sebagian pekerjaan terperincinya
telah kita lakukan dalam 262M.   Unsur dasar berikutnya adalah lemma
di bawah ini.
}%end of comment

\leader{263C}{Lemma} Misalkan $T$ suatu matriks real $r\times r$;  tetapkan
$J=|\det T|$.
Maka untuk setiap $\epsilon>0$ terdapat $\zeta=\zeta(T,\epsilon)>0$ sedemikian
sehingga

(i) $|\det S-\det T|\le\epsilon$ kapan pun $S$ adalah matriks $r\times r$
dan $\|S-T\|\le\zeta$;

(ii) kapan pun $D\subseteq\BbbR^r$ merupakan himpunan terbatas dan
$\phi:D\to\BbbR^r$ adalah fungsi sedemikian sehingga
$\|\phi(x)-\phi(y)-T(x-y)\|\le\zeta\|x-y\|$ untuk semua $x$, $y\in D$, maka
$|\mu^*\phi[D]-J\mu^*D|\le\epsilon\mu^*D$.

\proof{{\bf (a)} Tentu saja (i) adalah bagian yang mudah.   Karena
$\det S$ merupakan fungsi kontinu dari koefisien-koefisien $S$, dan
koefisien-koefisien $S$ harus dekat dengan koefisien-koefisien $T$ jika $\|S-T\|$
kecil (262Hb), tentu terdapat $\zeta_0>0$ sedemikian sehingga
$|\det S-\det T|\le\epsilon$ kapan pun $\|S-T\|\le\zeta_0$.

\medskip

{\bf (b)(i)} Tuliskan $B=B(\tbf{0},1)$ untuk bola satuan dalam $\BbbR^r$, dan
perhatikan $T[B]$.      Kita tahu bahwa $\mu T[B]=J\mu B$ (263A).    Misalkan
$G\supseteq T[B]$ suatu himpunan terbuka sedemikian sehingga $\mu G\le (J+\epsilon)\mu B$
(134Fa sekali lagi).   Karena $B$ kompak (2A2F sekali lagi), $T[B]$ juga kompak, sehingga
terdapat $\zeta_1>0$ sedemikian sehingga $T[B]+\zeta_1B\subseteq G$ (2A2Ed).
Ini berarti bahwa $\mu^*(T[B]+\zeta_1B)\le(J+\epsilon)\mu B$.

\medskip

\quad{\bf (ii)}
Sekarang andaikan bahwa $D\subseteq\BbbR^r$ merupakan himpunan terbatas, dan
$\phi:D\to\BbbR^r$ adalah fungsi sedemikian sehingga
$\|\phi(x)-\phi(y)-T(x-y)\|\le\zeta_1\|x-y\|$ untuk semua $x$, $y\in D$.
Maka jika $x\in D$ dan $\delta>0$,

\Centerline{$\phi[D\cap B(x,\delta)]\subseteq \phi(x) + \delta T[B] +
\delta\zeta_1B$,}

\noindent karena jika $y\in D\cap B(x,\delta)$ maka $T(y-x)\in\delta
T[B]$ dan

$$\eqalign{\phi(y)
&= \phi(x) + T(y-x) + (\phi(y)-\phi(x)-T(y-x))\cr
&\in\phi(x) + \delta T[B]+\zeta_1\|y-x\|B
\subseteq\phi(x)+\delta T[B]+\zeta_1\delta B.\cr}$$

\noindent Dengan demikian

$$\eqalign{\mu^*\phi[D\cap B(x,\delta)]
&\le\mu^*(\delta T[B]+\delta\zeta_1B)
=\delta^r\mu^*(T[B]+\zeta_1B)\cr
&\le\delta^r(J+\epsilon)\mu B
=(J+\epsilon)\mu B(x,\delta).\cr}$$

Misalkan $\eta>0$.   Maka terdapat barisan bola $\sequencen{B_n}$ dalam
$\BbbR^r$ sedemikian sehingga $D\subseteq\bigcup_{n\in\Bbb N}B_n$,
$\sum_{n=0}^{\infty}\mu B_n\le\mu^*D+\eta$ dan jumlah ukuran
bola-bola $B_n$ yang pusatnya tidak berada dalam $D$ paling besar $\eta$ (261F).
Misalkan $K$ adalah himpunan semua $n$ sedemikian sehingga pusat $B_n$ berada dalam
$D$.   Maka $\mu^*\phi[D\cap B_n]\le(J+\epsilon)\mu B_n$ untuk setiap
$n\in K$.    Juga, tentu saja, $\phi$ bersifat Lipschitz dengan konstanta
$(\|T\|+\zeta_1)$, sehingga $\mu^*\phi[D\cap B_n]\le(\|T\|+\zeta_1)^r\mu B_n$
untuk $n\in\Bbb N\setminus K$ (262D).   Sekarang

$$\eqalign{\mu^*\phi[D]
&\le\sum_{n=0}^{\infty}\mu^*\phi[D\cap B_n]\cr
&\le\sum_{n\in K}(J+\epsilon)\mu B_n
   +\sum_{n\in\Bbb N\setminus K}(\|T\|+\zeta_1)^r\mu B_n\cr
&\le(J+\epsilon)(\mu^*D+\eta)+\eta(\|T\|+\zeta_1)^r.\cr}$$

\noindent Karena $\eta$ sebarang,

\Centerline{$\mu^*\phi[D]\le (J+\epsilon)\mu^*D$.}

\medskip

{\bf (c)} Jika $J=0$, kita dapat berhenti di sini dengan menetapkan
$\zeta=\min(\zeta_0,\zeta_1)$;  sebab saat itu tentu berlaku $|\det S-\det
T|\le\epsilon$ kapan pun $\|S-T\|\le\zeta$, sedangkan jika $\phi:D\to\Bbb
R^r$ sedemikian sehingga $\|\phi(x)-\phi(y)-T(x-y)\|\le\zeta\|x-y\|$ untuk semua
$x$, $y\in D$, maka

\Centerline{$|\mu^*\phi[D]-J\mu^*D|=\mu^*\phi[D]\le\epsilon\mu^*D$.}

\noindent Jika $J\ne 0$, masih ada yang harus dilakukan.   Karena $T$ mempunyai
determinan tak nol, ia mempunyai invers $T^{-1}$, dan $|\det T^{-1}|=J^{-1}$.
Seperti dalam (b-i) di atas, terdapat $\zeta_2>0$ sedemikian sehingga
$\mu^*(T^{-1}[B]+\zeta_2B)\le(J^{-1}+\epsilon')\mu B$, dengan
$\epsilon'=\epsilon/(J(J+\epsilon))$.   Dengan mengulangi (b), kita melihat bahwa jika
$C\subseteq\Bbb R^r$ terbatas dan $\psi:C\to\BbbR^r$ sedemikian sehingga
$\|\psi(u)-\psi(v)-T^{-1}(u-v)\|\le\zeta_2\|u-v\|$ untuk semua $u$,
$v\in C$, maka $\mu^*\psi[C]\le(J^{-1}+\epsilon')\mu^*C$.

Sekarang andaikan bahwa
$D\subseteq\BbbR^r$ terbatas dan $\phi:D\to\BbbR^r$ sedemikian sehingga
$\|\phi(x)-\phi(y)-T(x-y)\|\le\zeta_2' \|x-y\|$ untuk semua $x$, $y\in D$,
dengan $\zeta_2'=\min(\zeta_2,\|T^{-1}\|)/2\|T^{-1}\|^2>0$.   Maka

\Centerline{$\|T^{-1}(\phi(x)-\phi(y))-(x-y)\|
\le\|T^{-1}\|\zeta_2'\|x-y\|\le\Bover12\|x-y\|$}

\noindent untuk semua $x$, $y\in D$, sehingga $\phi$ harus injektif;  tetapkan
$C=\phi[D]$ dan $\psi=\phi^{-1}:C\to D$.   Perhatikan bahwa $C$ terbatas,
karena

\Centerline{$\|\phi(x)-\phi(y)\|\le(\|T\|+\zeta_2')\|x-y\|$}

\noindent kapan pun $x$, $y\in D$.   Juga

\Centerline{$\|T^{-1}(u-v)-(\psi(u)-\psi(v))\|
\le\|T^{-1}\|\zeta_2'\|\psi(u)-\psi(v)\|\le\Bover12\|\psi(u)-\psi(v)\|$}

\noindent untuk semua $u$, $v\in C$.   Namun ini berarti bahwa

\Centerline{$\|\psi(u)-\psi(v)\|-\|T^{-1}\|\|u-v\|
\le\Bover12\|\psi(u)-\psi(v)\|$}

\noindent dan $\|\psi(u)-\psi(v)\|\le 2\|T^{-1}\|\|u-v\|$ untuk semua $u$,
$v\in C$, sehingga

\Centerline{$\|\psi(u)-\psi(v)-T^{-1}(u-v)\|
\le 2\zeta_2'\|T^{-1}\|^2\|u-v\|\le\zeta_2\|u-v\|$}

\noindent untuk semua $u$, $v\in C$.

Akibatnya

\Centerline{$\mu^*D=\mu^*\psi[C]\le(J^{-1}+\epsilon')\mu^*C
=(J^{-1}+\epsilon')\mu^*\phi[D]$,}

\noindent dan

\Centerline{$J\mu^*D\le (1+J\epsilon')\mu^*\phi[D]$.}

\medskip

{\bf (d)} Jadi jika kita menetapkan $\zeta=\min(\zeta_0,\zeta_1,\zeta_2')>0$, dan jika
$D\subseteq\BbbR^r$, $\phi:D\to\BbbR^r$ sedemikian sehingga $D$ terbatas
dan
$\|\phi(x)-\phi(y)-T(x-y)\|\le\zeta\|x-y\|$ untuk semua $x$, $y\in D$, kita
akan mempunyai

\Centerline{$\mu^*\phi[D]\le(J+\epsilon)\mu^*D$,}

\Centerline{$\mu^*\phi[D]\ge J\mu^*D-J\epsilon'\mu^*\phi[D]
\ge J\mu^*D-J\epsilon'(J+\epsilon)\mu^*D
= J\mu^*D-\epsilon\mu^*D$,}

\noindent sehingga kita memperoleh rumus yang diperlukan

\Centerline{$|\mu^*\phi[D]-J\mu^*D|\le\epsilon\mu^*D$.}
}%end of proof of 263C

\leader{263D}{}\cmmnt{ Kita siap untuk teorema ini.

\medskip

\noindent}{\bf Teorema} Misalkan $D\subseteq\BbbR^r$ sebarang himpunan, dan
$\phi:D\to\BbbR^r$ suatu fungsi yang terdiferensialkan relatif terhadap domainnya
di setiap titik dalam $D$.   Untuk setiap $x\in D$ misalkan $T(x)$ suatu turunan
$\phi$ relatif terhadap $D$ di $x$, dan tetapkan $J(x)=|\det T(x)|$.   Maka

\quad (i) $J:D\to\coint{0,\infty}$ merupakan fungsi terukur,

\quad (ii) $\mu^*\phi[D]\le\int_DJ\,d\mu$,

\noindent dengan membolehkan $\infty$ sebagai nilai integral.  Jika $D$
terukur, maka

\quad (iii) $\phi[D]$ terukur.

\noindent Jika $D$ terukur dan $\phi$ injektif, maka

\quad (iv) $\mu\phi[D]=\int_DJ\,d\mu$,

\quad (v) untuk setiap fungsi bernilai real $g$ yang didefinisikan pada suatu
himpunan bagian dari $\phi[D]$,

\Centerline{$\int_{\phi[D]}g\,d\mu=\int_{D}J\times g\phi\,d\mu$}

\noindent jika salah satu integral terdefinisi dalam $[-\infty,\infty]$, dengan syarat
kita menafsirkan $J(x)g(\phi(x))$ sebagai nol ketika $J(x)=0$ dan
$g(\phi(x))$ tidak terdefinisi.

\proof{{\bf (a)} Untuk melihat bahwa $J$ terukur, gunakan 262P;  fungsi
$T\mapsto|\det T|$ merupakan fungsi kontinu dari koefisien-koefisien $T$,
dan koefisien-koefisien $T(x)$ merupakan fungsi-fungsi terukur dari $x$, berdasarkan 262P,
sehingga $x\mapsto|\det T(x)|$ terukur (121K).   Kita juga tahu bahwa jika
$D$ terukur, $\phi[D]$ akan terukur, berdasarkan 262Ob.   Dengan demikian (i)
dan (iii) selesai.

\medskip

{\bf (b)} Untuk sementara, andaikan
bahwa $D$ terbatas, dan tetapkan $\epsilon>0$.    Untuk matriks-matriks $r\times r$
$T$, ambil $\zeta(T,\epsilon)>0$ seperti dalam 263C.   Ambil $\sequencen{D_n}$,
$\sequencen{T_n}$ seperti dalam 262M, sehingga $\sequencen{D_n}$ merupakan
penutup saling lepas dari $D$ oleh himpunan-himpunan yang terukur relatif dalam $D$, dan setiap
$T_n$ merupakan matriks $r\times r$ sedemikian sehingga

\Centerline{$\|T(x)-T_n\|\le\zeta(T_n,\epsilon)$ kapan pun $x\in D_n$,}

\Centerline{$\|\phi(x)-\phi(y)-T_n(x-y)\|\le\zeta(T_n,\epsilon)\|x-y\|$
untuk semua $x$, $y\in D_n$.}

\noindent Maka, dengan menetapkan $J_n=|\det T_n|$, kita mempunyai

\Centerline{$|J(x)-J_n|\le\epsilon$ untuk setiap $x\in D_n$,}

\Centerline{$|\mu^*\phi[D_n]-J_n\mu^*D_n|\le\epsilon\mu^*D_n$,}

\noindent berdasarkan pemilihan $\zeta(T_n,\epsilon)$.   Jadi kita mempunyai

\Centerline{$\int_DJ\,d\mu
\le\sum_{n=0}^{\infty}J_n\mu^*D_n+\epsilon\mu^*D
\le\int_DJ\,d\mu+2\epsilon\mu^*D$;}

\noindent di sini saya menggunakan fakta bahwa semua $D_n$ terukur relatif
dalam $D$, sehingga, khususnya,
$\mu^*D=\sum_{n=0}^{\infty}\mu^*D_n$.   Selanjutnya,

\Centerline{$\mu^*\phi[D]\le\sum_{n=0}^{\infty}\mu^*\phi[D_n]
\le\sum_{n=0}^{\infty}J_n\mu^*D_n+\epsilon\mu^*D$.}

\noindent Dengan menggabungkannya,

\Centerline{$\mu^*\phi[D]\le\int_DJ\,d\mu+2\epsilon\mu^*D$.}

Jika $D$ terukur dan $\phi$ injektif, maka semua $D_n$ merupakan
himpunan-himpunan bagian terukur dari $\BbbR^r$, sehingga semua $\phi[D_n]$ terukur,
dan himpunan-himpunan itu juga saling lepas.   Dengan demikian

\Centerline{$\int_DJ\,d\mu
\le\sum_{n=0}^{\infty}J_n\mu D_n+\epsilon\mu D
\le\sum_{n=0}^{\infty}(\mu\phi[D_n]+\epsilon\mu D_n)+\epsilon\mu D
=\mu\phi[D]+2\epsilon\mu D$.}

Karena $\epsilon$ sebarang, kita memperoleh

\Centerline{$\mu^*\phi[D]\le\int_DJ\,d\mu$,}

\noindent dan jika $D$ terukur dan $\phi$ injektif,

\Centerline{$\int_DJ\,d\mu\le\mu\phi[D]$;}

\noindent dengan demikian kita memperoleh (ii) dan (iv), dengan asumsi bahwa $D$
terbatas.

\medskip

{\bf (c)} Untuk himpunan umum $D$, tetapkan $B_k=B(\tbf{0},k)$;  maka

\Centerline{$\mu^*\phi[D]=\lim_{k\to\infty}\mu^*\phi[D\cap B_k]
\le\lim_{k\to\infty}\int_{D\cap B_k}J\,d\mu
=\int_DJ\,d\mu$,}

\noindent dengan kesamaan jika $\phi$ injektif dan $D$ terukur.

\medskip

{\bf (d)} Untuk bagian (v), saya hendak menunjukkan bahwa
hipotesis-hipotesis 235J terpenuhi, dengan mengambil $X=D$ dan $Y=\phi[D]$.   \Prf\
Tetapkan $G=\{x:x\in D,\,J(x)>0\}$.

\medskip

\quad\grheada\ Jika $F\subseteq\phi[D]$ terukur, maka terdapat
himpunan-himpunan Borel $F_1$, $F_2$ sedemikian sehingga
$F_1\subseteq F\subseteq F_2$ dan $\mu(F_2\setminus F_1)=0$.   Tetapkan
$E_j=\phi^{-1}[F_j]$ untuk setiap $j$, sehingga
$E_1\subseteq \phi^{-1}[F]\subseteq E_2$, dan kedua himpunan $E_j$
terukur,
karena $\phi$ dan $\dom\phi$ terukur.   Sekarang, dengan menerapkan (iv) pada
$\phi\restr E_j$,

\Centerline{$\int_{E_j}J\,d\mu=\mu\phi[E_j]=\mu(F_j\cap\phi[D])=\mu F$}

\noindent untuk kedua $j$, sehingga $\int_{E_2\setminus E_1}J\,d\mu=0$ dan $J=0$
hampir di mana-mana pada $E_2\setminus E_1$.   Dengan demikian
$J\times\chi(\phi^{-1}[F])\eae J\times\chi E_1$, dan
$\int J\times\chi(\phi^{-1}[F])d\mu$ ada dan sama dengan
$\int_{E_1}J\,d\mu=\mu F$.   Pada saat yang sama,
$(\phi^{-1}[F]\cap G)\symmdiff(E_1\cap G)$ terabaikan, sehingga
$\phi^{-1}[F]\cap G$ terukur.

\medskip

\quad\grheadb\ Jika $F\subseteq\phi[D]$ dan $G\cap\phi^{-1}[F]$
terukur, maka kita tahu bahwa $\mu\phi[D\setminus G]=\int_{D\setminus
G}J=0$ (berdasarkan (iv) yang diterapkan pada $\phi\restr D\setminus G$),
sehingga $F\setminus\phi[G]$ harus terabaikan;  sedangkan
$F\cap\phi[G]=\phi[G\cap\phi^{-1}[F]]$ juga terukur, berdasarkan (iii).
Dengan demikian $F$ terukur kapan pun $G\cap\phi^{-1}[F]$
terukur.

\medskip

Jadi semua hipotesis 235J terpenuhi.\ \QeD\   Sekarang (v) dapat
langsung dibaca dari kesimpulan 235J.
}%end of proof of 263D

\cmmnt{
\leader{263E}{Catatan (a)} Ini adalah suatu versi hasil klasik mengenai perubahan
variabel dalam integral berdimensi banyak.   Yang di sini saya sebut $J(x)$ adalah
{\bf Jacobian} dari $\phi$ di $x$;  ia menggambarkan perubahan volume
objek-objek di dekat $x$, mengikuti kaidah yang telah ditetapkan dalam 263A bagi
fungsi-fungsi dengan turunan konstan.   Gagasan
pembuktiannya juga merupakan gagasan klasik:  memecah himpunan $D$ menjadi
bagian-bagian $D_m$ yang cukup kecil agar kita dapat menghampiri $\phi$ dengan operator
afin $y\mapsto \phi(x) + T_{m}(y-x)$ pada masing-masing bagian.   Potensi
ketakteraturan himpunan $D$, yang dalam teorema ini boleh berupa sebarang himpunan,
diimbangi oleh kebebasan yang bersesuaian dalam
memilih himpunan-himpunan $D_m$.   Sesungguhnya terdapat dekomposisi lebih lanjut atas
himpunan-himpunan $D_m$ yang tersembunyi dalam bagian (b-ii) pembuktian 263C;  setiap $D_m$
pada hakikatnya ditutupi oleh suatu keluarga bola yang saling lepas, dan ukuran
citranya dapat kita taksir dengan ketelitian yang memadai.   Selalu ada bahaya
munculnya himpunan pengecualian yang terabaikan, dan kita memerlukan pertaksamaan kasar
dalam pembuktian 262D untuk menanganinya.

\header{263Eb}{\bf (b)} Di sepanjang pekerjaan dalam bab ini, dari 261B
sampai 263D, saya telah memilih
bola-bola $B(x,\delta)$ sebagai bentuk dasar yang digunakan.   Kiranya jelas
bahwa sebenarnya sebarang bentuk yang wajar juga akan sama baiknya.    Khususnya,
bentuk-bentuk yang disebut bola

\Centerline{$B_1(x,\delta)
=\{y:\sum_{i=1}^r|\eta_i-\xi_i|\le\delta\}$,\quad
$B_{\infty}(x,\delta)=\{y:|\eta_i-\xi_i|\le\delta\Forall
i\}$}

\noindent akan berfungsi dengan sempurna.   Ada banyak alternatif.   Kita
dapat menggunakan himpunan berbentuk $C(x,k)$, untuk $x\in\BbbR^r$ dan
$k\in\Bbb N$, yang didefinisikan sebagai kubus setengah-terbuka berbentuk
$\coint{2^{-k}z,2^{-k}(z+\tbf{1})}$ dengan $z\in\Bbb Z^r$ yang memuat $x$,
sebagai gantinya;  atau bahkan $C'(x,\delta)=\coint{x,x+\delta\tbf{1}}$.   Dalam semua
kasus semacam ini kita mempunyai versi teorema kepadatan (261Yb-261Yc) yang
mendukung teori selanjutnya.


\header{263Ec}{\bf (c)} Saya telah menyajikan 263D sebagai teorema mengenai
fungsi-fungsi terdiferensialkan, karena itulah bentuk normal penggunaannya dalam
penerapan elementer.   Namun, pembuktiannya pada hakikatnya bergantung pada
fakta bahwa fungsi terdiferensialkan merupakan gabungan terhitung dari fungsi-fungsi Lipschitz,
dan 263D akan langsung mengikuti dari teorema yang sama
yang dibuktikan hanya untuk fungsi-fungsi Lipschitz.   Faktanya, teorema ini
berlaku bagi {\it sebarang} gabungan terhitung dari fungsi-fungsi Lipschitz,
karena fungsi Lipschitz terdiferensialkan hampir di mana-mana.   Untuk
pekerjaan yang lebih lanjut (lihat {\smc Federer 69} atau {\smc Evans \& Gariepy 92},
atau Bab 47 dalam Jilid 4), tampaknya jelas bahwa fungsi-fungsi Lipschitz
merupakan objek yang penting, sehingga saya nyatakan hasilnya secara lengkap.
}%end of comment

\leader{*263F}{Akibat} Misalkan $D\subseteq\BbbR^r$ sebarang himpunan dan
$\phi:D\to\BbbR^r$ suatu fungsi Lipschitz.   Misalkan $D_1$ adalah himpunan
titik-titik tempat $\phi$ mempunyai turunan relatif terhadap $D$, dan untuk setiap
$x\in D_1$ misalkan $T(x)$ turunan semacam itu, dengan $J(x)=|\det T(x)|$.
Maka

\quad (i) $D\setminus D_1$ terabaikan;

\quad (ii) $J:D_1\to\coint{0,\infty}$ terukur;

\quad (iii) $\mu^*\phi[D]\le\int_DJ(x)dx$.

\noindent Jika $D$ terukur, maka

\quad (iv) $\phi[D]$ terukur.

\noindent Jika $D$ terukur dan $\phi$ injektif, maka

\quad (v) $\mu\phi[D]=\int_DJ\,d\mu$,

\quad (vi) untuk setiap fungsi bernilai real $g$ yang didefinisikan pada suatu
himpunan bagian dari $\phi[D]$,

\Centerline{$\int_{\phi[D]}g\,d\mu=\int_{D}J\times g\phi\,d\mu$}

\noindent jika salah satu integral terdefinisi dalam $[-\infty,\infty]$, dengan syarat
kita menafsirkan $J(x)g(\phi(x))$ sebagai nol ketika $J(x)=0$ dan
$g(\phi(x))$ tidak terdefinisi.

\proof{ Ini sekarang cukup dengan menggabungkan 262Q dan 263D,
dengan sedikit bantuan dari 262D.   Gunakan 262Q untuk menunjukkan bahwa
$D\setminus D_1$ terabaikan, 262D untuk menunjukkan bahwa $\phi[D\setminus D_1]$
terabaikan, dan terapkan 263D pada $\phi\restr D_1$.
}%end of proof of 263F

\dvro{\leader{263G}{Proposisi}
$\int_{-\infty}^{\infty}e^{-t^2/2}dt=\sqrt{2\pi}$.}
{\leader{263G}{Koordinat polar pada bidang} Saya berikan suatu contoh elementer
dengan akibat yang berguna.   Definisikan $\phi:\BbbR^2\to\BbbR^2$
dengan menetapkan $\phi(\rho,\theta)=(\rho\cos\theta,\rho\sin\theta)$ untuk
$\rho$, $\theta\in\BbbR$.   Maka $\phi'(\rho,\theta)
=\Matrix{\cos\theta&-\rho\sin\theta\\ \sin\theta&\rho\cos\theta}$, sehingga
$J(\rho,\theta)=|\rho|$ untuk semua $\rho$, $\theta$.   Tentu saja $\phi$ tidak
injektif, tetapi jika kita membatasinya pada domain
$D=\{(0,0)\}\cup\{(\rho,\theta):\rho>0,\,-\pi<\theta\le\pi\}$ maka
$\phi\restr D$ adalah bijeksi antara $D$ dan $\BbbR^2$, dan

\Centerline{$\int g\,d\xi_1d\xi_2=\int_Dg(\phi(\rho,\theta))\rho\,d\rho
d\theta$}

\noindent untuk setiap fungsi bernilai real $g$ yang terintegralkan pada
$\BbbR^2$.

Andaikan, khususnya, bahwa kita menetapkan

\Centerline{$g(x)=e^{-\|x\|^2/2}
=e^{-\xi_1^2/2}e^{-\xi_2^2/2}$}

\noindent untuk $x=(\xi_1,\xi_2)\in\Bbb R^2$.   Maka

\Centerline{$\int g(x)dx=\int e^{-\xi_1^2/2}d\xi_1
\int e^{-\xi_2^2/2}d\xi_2$,}

\noindent seperti dalam 253D.   Dengan menetapkan $I=\int e^{-t^2/2}dt$, kita mempunyai
$\int g=I^2$.   (Untuk melihat bahwa $I$ terdefinisi dengan baik dalam $\Bbb R$, perhatikan
bahwa integrannya kontinu, sehingga terukur, dan bahwa

\Centerline{$\int_{-1}^1e^{-t^2/2}dt\le 2$,}

\Centerline{$\int_{-\infty}^{-1}e^{-t^2/2}dt
=\int_1^{\infty}e^{-t^2/2}dt
\le\int_1^{\infty}e^{-t/2}dt
=\lim_{a\to\infty}\int_1^ae^{-t/2}dt
=2e^{-1/2}$}

\noindent keduanya berhingga.)   Sekarang, dengan memperhatikan bentuk alternatif, kita
memperoleh

$$\eqalignno{I^2
&=\int g(x)dx
=\int_Dg(\rho\cos\theta,\rho\sin\theta)\rho\,d(\rho,\theta)\cr
&=\int_De^{-\rho^2/2}\rho\,d(\rho,\theta)
=\int_0^{\infty}\int_{-\pi}^{\pi}\rho e^{-\rho^2/2}d\theta d\rho\cr
\noalign{\noindent (dengan mengabaikan titik $(0,0)$, yang berukuran nol)}
&=\int_0^{\infty}2\pi\rho e^{-\rho^2/2}d\rho
=2\pi\lim_{a\to\infty}\int_0^a\rho e^{-\rho^2/2}d\rho\cr
&=2\pi\lim_{a\to\infty}(-e^{-a^2/2}+1)
=2\pi.\cr}$$

\noindent Akibatnya

\Centerline{$\int_{-\infty}^{\infty}e^{-t^2/2}dt=I=\sqrt{2\pi}$,}

\noindent yang merupakan salah satu dari banyak fakta yang seharusnya diketahui setiap
matematikawan, dan khususnya sangat penting bagi Bab 27 di bawah.
}%end of dvro

\leader{263H}{Akibat} Jika $k\in\Bbb N$ ganjil,

\Centerline{$\int_{-\infty}^{\infty}x^ke^{-x^2/2}dx=0$;}

\noindent jika $k=2l\in\Bbb N$ genap, maka

\Centerline{$\int_{-\infty}^{\infty}x^ke^{-x^2/2}dx
=\Bover{(2l)!}{2^ll!}\sqrt{2\pi}$.}

\proof{{\bf (a)} Untuk melihat bahwa semua integral terdefinisi dengan baik dan
berhingga, perhatikan bahwa
$\lim_{x\to\pm\infty}x^ke^{-x^2/4}=0$, sehingga
$M_k=\sup_{x\in\Bbb R}|x^ke^{-x^2/4}|$ berhingga, dan

\Centerline{$\int_{-\infty}^{\infty}|x^ke^{-x^2/2}|dx
\le M_k\int_{-\infty}^{\infty}e^{-x^2/4}dx<\infty$.}

\medskip

{\bf (b)} Jika $k$ ganjil, maka dengan substitusi $y=-x$ kita memperoleh

\Centerline{$\int_{-\infty}^{\infty}x^ke^{-x^2/2}dx
=-\int_{-\infty}^{\infty}y^ke^{-y^2/2}dy$,}

\noindent sehingga kedua integral harus nol.

\medskip

{\bf (c)} Untuk $k$ genap, lanjutkan dengan induksi.   Tetapkan
$I_l=\int_{-\infty}^{\infty}x^{2l}e^{-x^2/2}dx$.
$I_0=\sqrt{2\pi}=\bover{0!}{2^00!}\sqrt{2\pi}$ berdasarkan 263G.   Untuk
langkah induksi menuju $l+1\ge 1$, lakukan integrasi parsial untuk melihat bahwa

\Centerline{$\int_{-a}^ax^{2l+1}\cdot xe^{-x^2/2}dx
=-a^{2l+1}e^{-a^2/2}+(-a)^{2l+1}e^{-a^2/2}
  +\int_{-a}^a(2l+1)x^{2l}e^{-x^2/2}dx$}

\noindent untuk setiap $a\ge 0$.   Dengan membiarkan $a\to\infty$,

\Centerline{$I_{l+1}=(2l+1)I_l$.}

\noindent Karena

\Centerline{$\Bover{(2(l+1))!}{2^{l+1}(l+1)!}\sqrt{2\pi}
=(2l+1)\Bover{(2l)!}{2^ll!}\sqrt{2\pi}$,}

\noindent induksi berlanjut.
}%end of proof of 263H

\leader{263I}{}\cmmnt{ Berikut ini adalah versi 263D untuk transformasi
yang tidak injektif.

\medskip

\noindent}{\bf Teorema}\dvAnew{2013}
Misalkan $D\subseteq\BbbR^r$ suatu himpunan terukur, dan
$\phi:D\to\BbbR^r$ suatu fungsi yang terdiferensialkan relatif terhadap domainnya di
setiap titik dalam $D$.   Untuk setiap $x\in D$ misalkan $T(x)$ suatu turunan
$\phi$ relatif terhadap $D$ di $x$, dan tetapkan $J(x)=|\det T(x)|$.

(a) Misalkan $\nu$ adalah ukuran cacah pada $\BbbR^r$.   Maka
$\int_{\BbbR^r}\nu(\phi^{-1}[\{y\}])dy$ dan
$\int_DJ\,d\mu$ terdefinisi dalam $[0,\infty]$ dan bernilai sama.

(b) Misalkan $g$ suatu fungsi bernilai real yang didefinisikan pada suatu himpunan
bagian dari $\phi[D]$
sedemikian sehingga $\int_Dg(\phi(x))\det T(x)dx$ terdefinisi dalam $\Bbb R$,
dengan menafsirkan $g(\phi(x))\det T(x)$ sebagai nol ketika $\det T(x)=0$ dan
$g(\phi(x))$ tidak terdefinisi.   Tetapkan

\Centerline{$C=\{y:y\in\phi[D]$, $\phi^{-1}[\{y\}]$ berhingga$\}$,
\quad$R(y)=\sum_{x\in\phi^{-1}[\{y\}]}\sgn\det T(x)$}

\noindent untuk $y\in C$, dengan $\sgn(0)=0$ dan
$\sgn(\alpha)=\Bover{\alpha}{|\alpha|}$ untuk $\alpha$ tak nol.
Jika kita menafsirkan $g(y)R(y)$ sebagai nol ketika $g(y)=0$ dan $R(y)$ tidak terdefinisi,
maka $\int_{\phi[D]}g\times R\,d\mu$ terdefinisi dan sama dengan
$\int_Dg(\phi(x))\det T(x)dx$.

\proof{{\bf (a)} Berdasarkan 263D(i), $J$ terukur, sehingga
$\int_DJ\,d\mu$ terdefinisi dalam $[0,\infty]$ dan
$D_0=\{x:x\in D$, $J(x)=0\}$ terukur.   Dengan menerapkan 263D(ii) pada
$\phi\restr D_0$, kita melihat bahwa $\phi[D_0]$ terabaikan.

Dengan menerapkan 262M pada $\phi\restr D\setminus D_0$,
himpunan $A$ dari matriks-matriks $r\times r$ nonsingular dan
$\zeta(S)=\bover1{2\|S^{-1}\|}$ untuk $S\in A$,
kita memperoleh partisi $\sequencen{E_n}$ atas $D\setminus D_0$ menjadi himpunan-himpunan
terukur dan barisan $\sequencen{T_n}$ dalam $A$ sedemikian sehingga

\Centerline{$\|\phi(x)-\phi(y)-T_n(x-y)\|\le\Bover1{2\|T_n^{-1}\|}\|x-y\|$,
\quad$\|T(x)-T_n\|\le\Bover1{2\|T_n^{-1}\|}$}

\noindent kapan pun
$n\in\Bbb N$ dan $x$, $y\in E_n$.   Dalam hal ini, untuk $x$, $y\in E_n$,

\Centerline{$\phi(x)=\phi(y)
\,\Longrightarrow\,\|x-y\|\le\|T_n^{-1}\|\|T_n(x-y)\|
\le\Bover12\|x-y\|
\,\Longrightarrow\, x=y$,}

\noindent sehingga $\phi\restr E_n$ injektif, untuk setiap $n$.
Akibatnya
$\#(\phi^{-1}[\{y\}])=\#(\{n:y\in\phi[E_n]\})$ untuk
$y\in\phi[D]\setminus\phi[D_0]$, dan

$$\eqalignno{\int_{\phi[D]}\nu(\phi^{-1}[\{y\}])dy
&=\int_{\phi[D]\setminus\phi[D_0]}\nu(\phi^{-1}[\{y\}])dy
=\sum_{n=0}^{\infty}\mu(\phi[E_n]\setminus\phi[D_0])\cr
\displaycause{dengan menerapkan 263D(iii) pada $\phi\restr E_n$, kita tahu bahwa
$\phi[E_n]$ terukur untuk setiap $n$}
&=\sum_{n=0}^{\infty}\mu\phi[E_n]
=\sum_{n=0}^{\infty}\int_{E_n}J\cr
\displaycause{dengan menerapkan 263D(iv) pada $\phi\restr E_n$}
&=\int_{D\setminus D_0}J
=\int_DJ,\cr}$$

\noindent dengan setiap jumlah atau integral terdefinisi dalam $[0,\infty]$ karena
yang berikutnya demikian.

\medskip

{\bf (b)(i)} Dengan menetapkan $D'=\phi^{-1}[\dom g]$, kita melihat bahwa
$D\setminus(D_0\cup D')$ terabaikan, sehingga (dengan menggunakan 263D(ii) lagi)

\Centerline{$\phi[D]\setminus\dom g=\phi[D]\setminus\phi[D']
\subseteq\phi[D_0]\cup\phi[D\setminus(D_0\cup D')]$}

\noindent terabaikan.   Selanjutnya, jika kita menetapkan
$C'=C\cup g^{-1}[\{0\}]$, $\phi[D]\setminus C'$ terabaikan.
\Prf\ Untuk setiap $m\in\Bbb N$, tetapkan
$F_m=\{y:y\in\dom g$, $|g(y)|\ge 2^{-m}\}$.   Maka

\Centerline{$\phi^{-1}[F_m]\setminus D_0=\{x:x\in D$, $J(x)\ne 0$,
   $g(\phi(x))$ terdefinisi dan $|g(\phi(x))|\ge 2^{-m}\}$}

\noindent terukur (karena $J\times g\phi$ terukur) dan
(dengan menerapkan (a) pada $\phi\restr\phi^{-1}[F_m]$)

\Centerline{$\int_{F_m}\nu(\phi^{-1}[\{y\}])dy
=\int_{\phi^{-1}[F_m]}J\,d\mu
\le 2^m\int_{\phi^{-1}[F_m]}|J\times g(\phi)|\,d\mu$}

\noindent berhingga.   Namun ini berarti bahwa $\nu(\phi^{-1}[\{y\}])$ harus
berhingga untuk hampir setiap $y\in F_m$, yakni, bahwa $F_m\setminus C$
terabaikan.   Karena $m$ sebarang, $\dom g\setminus C'$ dan
$\phi[D]\setminus C'$ terabaikan.\ \Qed

\medskip

\quad{\bf (ii)} Dengan mengambil $\sequencen{E_n}$ dan $\sequencen{T_n}$ seperti dalam (a),
tetapkan $\epsilon_n=\sgn\det T_n\in\{-1,1\}$ untuk setiap $n$.   Maka
$\sgn\det T(x)=\epsilon_n$ kapan pun $n\in\Bbb N$ dan $x\in E_n$.   \Prf\
Untuk setiap $\alpha\in[0,1]$,

\Centerline{$\|(\alpha T(x)+(1-\alpha)T_n)-T_n\|
\le\|T(x)-T_n\|\le\Bover1{2\|T_n^{-1}\|}$,}

\noindent sehingga (dengan menggunakan 2A4Fd)
$\|T_n^{-1}(\alpha T(x)+(1-\alpha)T_n)-I_r\|\le\Bover12$,
dengan $I_r$ matriks identitas $r\times r$,
dan $\alpha T(x)+(1-\alpha)T_n$ nonsingular
(sebab jika $(\alpha T(x)+(1-\alpha)T_n)z=0$, maka
$\|z\|\le\bover12\|z\|$).   Dengan demikian $\det(\alpha T(x)+(1-\alpha)T_n)$
tak nol untuk $0\le\alpha\le 1$.   Namun karena
$\alpha\mapsto\det(\alpha T(x)+(1-\alpha)T_n)$ kontinu,
$\epsilon_n=\sgn\det T_n=\sgn\det T(x)$.\ \Qed

\medskip

\quad{\bf (iii)} Sekarang

$$\eqalignno{\int_Dg(\phi(x))\det T(x)\,dx
&=\int_{D\setminus D_0}g(\phi(x))\det T(x)\,dx\cr
&=\sum_{n=0}^{\infty}\int_{E_n}g(\phi(x))\det T(x)\,dx\cr
\displaycause{dalam rangkaian rumus ini, setiap jumlah dan integral
terdefinisi dengan baik karena yang sebelumnya demikian}
&=\sum_{n=0}^{\infty}\epsilon_n\int_{E_n}g(\phi(x))J(x)\,dx
=\sum_{n=0}^{\infty}\epsilon_n\int_{\phi[E_n]}g\,dx\cr
&=\sum_{n=0}^{\infty}\epsilon_n\int_{\phi[E_n]\cap C'}g\,dx
=\sum_{n=0}^{\infty}\epsilon_n\int_{C'}g\times\chi(\phi[E_n])\,dx\cr}$$

\noindent (gunakan 131Fa, jika Anda mau, untuk langkah terakhir).
Karena kita juga mempunyai

\Centerline{$\infty
>\int_D|g(\phi(x))\det T(x)|\,dx
=\sum_{n=0}^{\infty}\int_{C'}|g|\times\chi(\phi[E_n])\,dx$}

\noindent (dengan melalui tahap-tahap yang sama menggunakan nilai mutlak
integran), kita mempunyai

$$\eqalignno{\int_Dg(\phi(x))\det T(x)\,dx
&=\sum_{n=0}^{\infty}\epsilon_n\int_{C'}g\times\chi(\phi[E_n])\,dx\cr
&=\int_{C'}g\times\sum_{n=0}^{\infty}\epsilon_n\chi(\phi[E_n])\,dx\cr
&=\int_{C'\setminus\phi[D_0]}
   g\times\sum_{n=0}^{\infty}\epsilon_n\chi(\phi[E_n])\,dx
=\int_{C'\setminus\phi[D_0]}g\times R\,d\mu\cr
\displaycause{karena jika $y\in C'\setminus\phi[D_0]$, $g(y)=0$ atau
$R(y)$ terdefinisi dan sama dengan
$\sum_{n=0}^{\infty}\epsilon_n\chi(\phi[E_n])(y)$}
&=\int_{\phi[D]}g\times R\,d\mu,\cr}$$

\noindent sebagaimana diklaim.
}%end of proof of 263I

\leader{263J}{Kasus \dvrocolon{berdimensi satu}}\cmmnt{ Pembatasan pada fungsi
injektif $\phi$ dalam 263D(v) tidak dapat dihindari dalam konteks hasil
tersebut.   Namun, dalam substitusi kalkulus elementer, pembatasan ini tidak
selalu perlu.   Dengan harapan memperjelas keadaan, di sini saya berikan
suatu hasil yang mencakup banyak kiat standar.

\medskip

\noindent}{\bf Proposisi}
Misalkan $I\subseteq\Bbb R$ suatu interval dengan lebih dari
satu titik, dan $\phi:I\to\Bbb R$ suatu fungsi yang kontinu mutlak
pada setiap subinterval tertutup terbatas dari $I$.   Tuliskan $u=\inf I$,
$u'=\sup I$ dalam
$[-\infty,\infty]$, dan andaikan bahwa $v=\lim_{x\downarrow u}\phi(x)$ dan
$v'=\lim_{x\uparrow u'}\phi(x)$ terdefinisi dalam $[-\infty,\infty]$.   Misalkan
$g$ suatu fungsi real sedemikian sehingga
$\int_Ig(\phi(x))\phi'(x)dx$ terdefinisi,
dengan ketentuan bahwa kita menafsirkan $g(\phi(x))\phi'(x)$ sebagai $0$ ketika
$\phi'(x)=0$ dan $g(\phi(x))$ tidak terdefinisi.   Maka
$\int_v^{v'}g$ terdefinisi dan sama dengan $\int_Ig(\phi(x))\phi'(x)dx$, dengan
di sini $\int_v^{v'}g$ ditafsirkan sebagai $-\int_{v'}^vg$ jika
$v'<v$.

\proof{{\bf (a)} $\phi$ terdiferensialkan hampir di mana-mana
pada $I$ dan $\phi[A]$ terabaikan untuk setiap himpunan terabaikan
$A\subseteq I$.   \Prf\ Kita dapat menyatakan $I$ sebagai gabungan suatu
barisan $\sequencen{I_n}$ dari interval-interval tertutup terbatas sedemikian sehingga
$\phi\restr I_n$ kontinu mutlak untuk setiap $n$.   Berdasarkan
225Cb dan 225G, $\phi\restr I_n$ terdiferensialkan hampir di mana-mana
pada $I_n$ dan $\phi[A]$ terabaikan untuk setiap himpunan terabaikan
$A\subseteq I_n$, untuk setiap $n$.   Jadi
$\phi$ terdiferensialkan hampir di mana-mana
pada $\bigcup_{n\in\Bbb N}I_n=I$ dan
$\phi[A]=\bigcup_{n\in\Bbb N}\phi[A\cap I_n]$ terabaikan untuk setiap himpunan
terabaikan $A\subseteq I$.\ \Qed

Karena $\phi\restr J$ kontinu untuk setiap interval tertutup
$J\subseteq I$, $\phi$ kontinu.   Berdasarkan Teorema Nilai Antara,
$\phi[I]$ merupakan interval yang memuat $\ooint{\min(v,v'),\max(v,v')}$.

\medskip

{\bf (b)} Misalkan $D\subseteq I$ adalah domain $\phi'$.   Untuk $x\in D$,
kita dapat memandang $\phi'(x)$ sebagai matriks $1\times 1$ dengan determinan
$\phi'(x)$.   Karena $I\setminus D$ terabaikan,
$\phi[I]\setminus\phi[D]\subseteq\phi[I\setminus D]$ terabaikan.
Sekarang $\int_Dg(\phi(x))\phi'(x)dx=\int_Ig(\phi(x))\phi'(x)dx$.
Dengan menerapkan 263I pada $\phi\restr D$ dan $g\restr\phi[D]$, kita melihat bahwa
$\int_{\phi[D]}g\times R$ terdefinisi dan sama dengan
$\int_Dg\phi\times\phi'$, dengan
$R(y)=\sum_{x\in D\cap\phi^{-1}[\{y\}]}\sgn\phi'(x)$ kapan pun
$y\in\phi[D]$ dan $D\cap\phi^{-1}[\{y\}]$ berhingga.

\medskip

{\bf (c)} (Langkah kunci.) Tetapkan $D_0=\{x:x\in D$, $\phi'(x)=0\}$.   Berdasarkan
263D(ii), yang diterapkan pada $\phi\restr D_0$, $\phi[D_0]$ terabaikan.
Tetapkan

\Centerline{$C=\{y:y\in\phi[D]\cap\dom g\setminus(\phi[D_0]\cup\{v,v'\})$,
$\phi^{-1}[\{y\}]$ berhingga, $g(y)\ne 0\}$.}

\noindent Jika $y\in C$ dan $K=\phi^{-1}[\{y\}]$, maka

$$\eqalign{R(y)
=\sum_{x\in K}\sgn\phi'(x)
&=1\text{ jika }v<y<v',\cr
&=-1\text{ jika }v'<y<v,\cr
&=0\text{ selain itu}.\cr}$$

\noindent\Prf\
Jika $J\subseteq I\setminus K$ suatu interval, $\phi(z)\ne y$ untuk
$z\in J$;  karena $\phi$ kontinu, Teorema Nilai Antara
memberi tahu kita bahwa $\sgn(\phi(z)-y)$ konstan pada $J$.
Juga $\phi'(x)\ne 0$ untuk setiap $x\in K$, karena $y\notin\phi[D_0]$.
Induksi sederhana
atas $\#(K\cap\ooint{-\infty,z})$ menunjukkan bahwa
$\sgn(\phi(z)-y)=\sgn(v-y)+2\sum_{x\in K,x<z}\sgn\phi'(x)$ untuk setiap
$z\in I\setminus K$;  dengan mengambil limit ketika $z\uparrow u'$,
$\sum_{x\in K}\sgn\phi'(x)=\bover12(\sgn(v'-y)-\sgn(v-y))$.
(Di sini mungkin kita harus menafsirkan $\sgn(\pm\infty)$ sebagai $\pm 1$
dengan cara yang wajar.)   Ini ternyata tepat seperti yang perlu kita ketahui.\ \Qed

\medskip

{\bf (d)} Jadi sekarang kita mempunyai

$$\eqalign{\int_Ig\phi\times\phi'
&=\int_Dg\phi\times\phi'
=\int_{\phi[D]}g\times R\cr
&=\int_{\phi[D]}g\times R\times\chi C
=\int_v^{v'}g\times\chi C
=\int_v^{v'}g\cr}$$

\noindent karena $\phi[D]\setminus(C\cup g^{-1}[\{0\}])$ terabaikan.
}%end of proof of 263J

\exercises{
\leader{263X}{Latihan dasar (a)}
%\spheader 263Xa
Misalkan $(X,\Sigma,\mu)$ sebarang ruang ukur, $f\in\eusm L^0(\mu)$ dan
$p\in\coint{1,\infty}$.   Tunjukkan bahwa $f\in\eusm L^p(\mu)$ jika dan hanya jika

\Centerline{$\gamma=p\int_0^{\infty}\alpha^{p-1}
\mu^*\{x:x\in\dom f,\,|f(x)|>\alpha\}d\alpha$}

\noindent berhingga, dan dalam hal ini $\|f\|_p=\gamma^{1/p}$.
\Hint{$\int|f|^p=\int_0^{\infty}\mu^*\{x:|f(x)|^p>\beta\}d\beta$, berdasarkan
252O;  sekarang substitusikan $\beta=\alpha^p$.}
%263D

\spheader 263Xb Misalkan $f$ suatu fungsi terintegralkan yang didefinisikan hampir
di mana-mana dalam $\BbbR^r$.   Tunjukkan bahwa jika $\alpha<r-1$ maka
$\sum_{n=1}^{\infty}n^{\alpha}|f(nx)|$ berhingga untuk hampir setiap
$x\in\BbbR^r$.   \Hint{Taksirlah
$\sum_{n=0}^{\infty}n^{\alpha}\int_B|f(nx)|dx$ untuk bola-bola $B$
yang berpusat di titik asal.}
%263D

\spheader 263Xc Misalkan $A\subseteq\ooint{0,1}$ suatu himpunan sedemikian sehingga
$\mu^*A=\mu^*([0,1]\setminus A)=1$, dengan $\mu$ ukuran Lebesgue pada
$\Bbb R$.   Tetapkan $D=A\cup\{-x:x\in\ooint{0,1}\setminus A\}
\subseteq[-1,1]$, dan tetapkan $\phi(x)=|x|$ untuk $x\in D$.   Tunjukkan bahwa $\phi$
injektif, bahwa $\phi$ terdiferensialkan relatif terhadap domainnya
pada setiap titik dalam $D$, dan bahwa $\mu^*\phi[D]<\int_D|\phi'(x)|dx$.
%263D

\spheader 263Xd Misalkan $\phi:D\to\BbbR^r$ suatu fungsi yang terdiferensialkan
relatif terhadap $D$ di setiap titik dari $D\subseteq\BbbR^r$, dan andaikan bahwa
untuk setiap
$x\in D$ terdapat turunan nonsingular $T(x)$ dari $\phi$ di $x$.
Tunjukkan bahwa $D$ dapat dinyatakan sebagai
$\bigcup_{k\in\Bbb N}D_k$ dengan $D_k=D\cap\overline{D}_k$ dan
$\phi\restr D_k$ injektif untuk setiap $k$.
%263D

\sqheader 263Xe(i) Tunjukkan bahwa untuk setiap himpunan terukur Lebesgue
$E\subseteq\Bbb R$ dan $t\in\Bbb R\setminus\{0\}$,
$\int_{tE}\bover1{|u|}du=\int_E\bover1{|u|}du$.   (ii) Untuk $t\in\Bbb R$,
$u\in\Bbb R\setminus\{0\}$ tetapkan $\phi(t,u)=(\bover{t}{u},u)$.   Tunjukkan bahwa
$\int_{\phi[E]}\bover1{|tu|}d(t,u)=\int_E\bover1{|tu|}d(t,u)$ untuk setiap
himpunan terukur Lebesgue $E\subseteq\BbbR^2$.
%263D

\sqheader 263Xf Definisikan $\phi:\BbbR^3\to\BbbR^3$ dengan
menetapkan

\Centerline{$\phi(\rho,\theta,\alpha)
=(\rho\sin\theta\sin\alpha,\rho\cos\theta\sin\alpha,\rho\cos\alpha)$.}

\noindent Tunjukkan bahwa $\det\phi'(\rho,\theta,\alpha)=\rho^2\sin\alpha$.
%263G

\spheader 263Xg Tunjukkan bahwa jika $k=2l+1$ ganjil, maka
$\int_0^{\infty}x^ke^{-x^2/2}dx=2^ll!$.   (Bandingkan 252Xi.)
%263H

\leader{263Y}{Latihan lanjutan (a)}
%\spheader 263Ya
Definisikan ukuran $\nu$ pada $\Bbb R$ dengan menetapkan
$\nu E=\int_E\bover1{|x|}dx$ untuk himpunan-himpunan terukur Lebesgue
$E\subseteq\Bbb R$.   Untuk $f$, $g\in\eusm L^1(\nu)$ tetapkan
$(f*g)(x)=\int f(\bover{x}{t})g(t)\nu(dt)$ kapan pun ini terdefinisi dalam
$\Bbb R$.   (i) Tunjukkan bahwa $f*g=g*f\in\eusm L^1(\nu)$.   (ii) Tunjukkan bahwa
$\int h(x)(f*g)(x)\nu(dx)=\int h(xy)f(x)g(y)\nu(dx)\nu(dy)$ untuk setiap
$h\in\eusm L^{\infty}(\nu)$.   (iii) Tunjukkan bahwa $f*(g*h)=(f*g)*h$ untuk
setiap $h\in\eusm L^1(\nu)$.   \Hint{263Xe.}
%263Xe 263D

\spheader 263Yb Misalkan $E\subseteq\BbbR^2$ suatu himpunan terukur sedemikian sehingga
$\limsup_{\alpha\to\infty}
\bover1{\alpha^2}\mu_2(E\cap B(\tbf{0},\alpha))>0$, dengan $\mu_2$ menyatakan
ukuran Lebesgue pada
$\BbbR^2$.   Tunjukkan bahwa terdapat $\theta\in\ocint{-\pi,\pi}$ sedemikian sehingga
$\mu_1E_{\theta}=\infty$, dengan
$E_{\theta}=\{\rho:\rho\ge 0,\,(\rho\cos\theta,\rho\sin\theta)\in E\}$.
\Hint{Tunjukkan bahwa $\mu_2(E\cap B(\tbf{0},\alpha))
\le\alpha\int_{-\pi}^{\pi}\mu_1E_{\theta}d\theta$.}
Generalisasikan ke dimensi yang lebih tinggi dan ke fungsi-fungsi selain $\chi E$.
%263G

\spheader 263Yc Misalkan $E\subseteq\BbbR^r$ suatu himpunan terukur,
dan $\phi:E\to\BbbR^r$ suatu fungsi yang terdiferensialkan relatif terhadap domainnya,
dengan turunan $T(x)$ di setiap titik $x$ dari $E$;  tetapkan
$J(x)=|\det T(x)|$.
Tunjukkan bahwa untuk setiap fungsi terintegralkan $g$ yang didefinisikan pada $\phi[E]$,

\Centerline{$\int g(y)\#(\phi^{-1}[\{y\}])dy
=\int_EJ(x)g(\phi(x))dx$.}
%263I

\spheader 263Yd Temukan pembuktian 263J berdasarkan gagasan-gagasan \S225.
\Hint{225Xe.}
%263J

\spheader 263Ye Misalkan $f:[a,b]\to\Bbb R$ suatu fungsi bervariasi terbatas,
dengan $a<b$ dalam $\Bbb R$, dengan dekomposisi Lebesgue
$f=f_p+f_{cs}+f_{ac}$ seperti dalam 226Cd;  misalkan $\mu$ ukuran Lebesgue pada
$\Bbb R$.   Tunjukkan bahwa pernyataan-pernyataan berikut ekuivalen:  (i) $f_{cs}$
konstan;  (ii) $\mu f[\,[c,d]\,]\le\int_c^d|f'|d\mu$ kapan pun
$a\le c\le d\le b$;   (iii)
$\mu^*f[A]\le\int_A|f'|d\mu$ untuk setiap $A\subseteq[a,b]$;   (iv) $f[A]$
terabaikan untuk setiap himpunan terabaikan $A\subseteq[a,b]$.   \Hint{Untuk
(iv)$\Rightarrow$(i), gabungkan 226Yd dan 263D(ii) untuk menunjukkan bahwa
$|f(d)-f(c)|\le\int_c^d|f'|d\mu+\Var_{[c,d]}f_p$ kapan pun
$a\le c\le d\le b$, dan karena itu
$\Var_{[a,b]}f\le\Var_{[a,b]}f_p+\Var_{[a,b]}f_{ac}$.}
%263D

\spheader 263Yf\dvAnew{2013} Andaikan $r=2$ dan
$\phi:\BbbR^2\to\BbbR^2$ terdiferensialkan secara kontinu dengan turunan nonsingular
$T$ di $\tbf{0}$.   (i) Tunjukkan bahwa terdapat $\epsilon>0$
sedemikian sehingga kapan pun
$\Gamma$ suatu lingkaran kecil berpusat di $\tbf{0}$ dengan jari-jari paling besar
$\epsilon$, maka
$\phi\restr\Gamma$ adalah homeomorfisme antara $\Gamma$ dan suatu kurva tertutup
sederhana yang mengelilingi $\tbf{0}$.   (ii) Tunjukkan bahwa jika $\det T>0$, maka
untuk lingkaran-lingkaran semacam itu $\phi(x)$ bergerak berlawanan arah jarum jam
mengelilingi $\phi[\Gamma]$ ketika $x$
bergerak berlawanan arah jarum jam mengelilingi $\Gamma$.   (iii) Apa yang terjadi jika
$\det T<0$?
}%end of exercises

\endnotes{
\Notesheader{263} Sekali lagi, ketika mendekati 263D, saya harus
memilih antara memberikan hasil yang mudah dipahami tetapi relatif lemah, dan
mencurahkan upaya tambahan untuk menyajikan teorema yang mendekati
batas alamiah dari apa yang dapat dicapai dengan konsep-konsep yang sedang
dikembangkan dalam jilid ini;  dan, seperti biasa, saya memilih bentuk yang lebih
kuat.   Ada tiga sumber kesulitan utama:  (i) fakta bahwa
kita menangani lebih dari satu dimensi;  (ii) fakta bahwa
kita menangani domain yang tidak teratur;   (iii) fakta bahwa
kita menangani fungsi terintegralkan yang sebarang.   Saya rasa saya tidak perlu
meminta maaf atas (iii) dalam buku mengenai teori ukuran.   Mengenai (ii),
memang benar bahwa penerapan utama hasil-hasil ini adalah pada
kasus-kasus ketika transformasi $\phi$ terdiferensialkan di mana-mana,
dengan turunan kontinu, dan himpunan $D$ mempunyai batas yang terabaikan;
dan dalam kasus-kasus ini tersedia penyederhanaan yang berarti --
terutama karena himpunan-himpunan $D_m$ dalam pembuktian 263D dapat diambil sebagai
kubus.   Meskipun demikian, menurut saya setiap bentuk hasil yang menggunakan
asumsi semacam itu sangat tidak memuaskan pada tingkat ini, karena merupakan
kompromi canggung antara gagasan yang alamiah bagi integral Riemann dan gagasan
yang alamiah bagi integral Lebesgue.   Mengenai (i), mungkin bahkan tepat untuk
menguraikan seluruh argumen untuk kasus $r=1$ sebelum
beralih ke kasus umum, seperti yang saya lakukan dalam \S\S114-115, karena
kasus berdimensi satu sudah penting dan menarik;  dan jika Anda
merasa pekerjaan di atas sulit -- dan memang sulit -- sementara minat langsung
Anda adalah integrasi berdimensi satu melalui substitusi, maka saya kira
bermanfaat bagi Anda untuk menyusun sendiri argumen $r=1$, sampai pada
pembuktian 263J.   Sesungguhnya perbedaan terbesar terdapat
dalam 263A, yang menjadi hampir sepele;  pekerjaan dalam 262M dan 263C
menjadi lebih mudah dibaca, karena semua matriks berubah menjadi skalar dan kita
dapat menghilangkan kata "determinan", tetapi menurut saya kita tidak dapat
meninggalkan satu pun gagasannya, setidaknya jika kita ingin memperoleh 263D
seperti yang dinyatakan.   (Namun lihat 263Yd.)

Dalam paragraf terakhir saya bersikeras bahwa suatu pembedaan dapat
dibuat antara "gagasan yang alamiah bagi integral Riemann dan gagasan
yang alamiah bagi integral Lebesgue".   Di sini kita mendekati pertanyaan yang
mendalam, seperti "untuk apa buku teori ukuran dibuat?", yang menurut saya
tidak dapat dijawab tanpa suatu acuan -- mungkin secara tidak sadar --
kepada pertanyaan "untuk apa matematika dibuat?".   Tentu saja saya
ingin menyajikan di sini beberapa teorema umum yang mengagumkan yang
muncul dalam teori Lebesgue.   Namun yang lebih penting daripada teorema tertentu
mana pun adalah gambaran umum tentang apa yang dapat dibuktikan dengan metode-metode ini.
Hakikat teori ukuran modern adalah bahwa kontinuitas tidak
penting, atau, jika Anda lebih suka, bahwa fungsi-fungsi terukur dalam arti tertentu
begitu mendekati kontinu sehingga kita tidak perlu menambahkan hipotesis
kontinuitas dalam teorema-teorema kita.   Dalam suatu arti, ini merupakan pembebasan
besar, dan integral Lebesgue kini menjadi integral standar.   Namun Anda tidak
boleh menganggap integral Riemann sudah usang.   Intuisi yang melandasinya --
misalnya, bahwa permukaan suatu benda pejal mempunyai volume nol -- tetap
sangat bernilai dalam konteks yang tepat, yang tentu mencakup kajian
fungsi terdiferensialkan dengan turunan kontinu.   Yang saya katakan
di sini ialah bahwa menurut saya kita dapat menggunakan intuisi tersebut sebaik-baiknya
jika kita mempertahankan pembagian, tentu saja yang lentur dan dapat ditembus,
antara gagasan kedua teori itu;  dan bahwa ketika memindahkan suatu teorema
dari satu sisi batas ke sisi lainnya, kita seharusnya melakukannya
secara sungguh-sungguh, dengan berusaha mengungkapkan seluruh kekuatan metode yang
kita gunakan.

Saya telah mengatakan bahwa perbedaan hakiki antara
kasus berdimensi satu dan berdimensi banyak terletak dalam 263A, tempat Jacobian
$J=|\det T|$ masuk ke dalam argumen.   Jika perangkat teknis yang diperlukan
untuk menangani sebarang himpunan terukur Lebesgue disisihkan, hal ini
merupakan perhitungan volume paralelipipedum $T[I]$
dengan $I$ interval $\coint{\tbf{0},\tbf{1}}$.   Saya menangani
hal ini dengan sedikit aljabar, dengan mengatakan bahwa hasilnya pada hakikatnya
jelas jika $T$ diagonal, sedangkan jika $T$ suatu isometri hasilnya mengikuti
dari fakta bahwa bola satuan tetap invarian;  dan aljabar digunakan untuk
menyatakan sebarang matriks sebagai produk matriks diagonal dan ortogonal.
Dari 261F juga jelas bahwa ukuran Lebesgue
harus invarian terhadap rotasi maupun terhadap translasi;
dengan kata lain, ia invarian di bawah semua isometri.   Cara lain memandang
hal ini akan muncul dalam bagian berikutnya.

Menurut saya pusat argumen untuk 263D terletak dalam lemma
263C.   Di sinilah kita mengubah hasil eksak untuk operator linear menjadi
hasil hampiran untuk fungsi yang hampir linear;  dan inti
keterdiferensialan ialah bahwa fungsi terdiferensialkan dihampiri dengan baik,
di dekat setiap titik domainnya, oleh operator linear.
Lemma ini melibatkan dua gagasan yang agak berbeda.   Untuk menunjukkan
bahwa $\mu^*\phi[D]\le(J+\epsilon)\mu^*D$, pertama-tama kita memperhatikan bola-bola
dan kemudian menggunakan teorema Vitali untuk melihat bahwa $D$ ditutupi
secara hemat oleh bola-bola, sehingga batas atas bagi $\mu^*\phi[D]$ dalam
bentuk jumlah $\sum_{B\in\Cal I_0}\mu^*\phi[D\cap B]$ sudah memadai.
Untuk memperoleh batas bawah, kita perlu membalik argumen dengan memperhatikan
$\psi=\phi^{-1}$, yang mengharuskan kita terlebih dahulu memeriksa bahwa $\phi$
dapat dibalik, lalu bahwa $\psi$ terhubung secara tepat dengan $T^{-1}$.
Saya telah menuliskan rumus eksak untuk $\epsilon'$, $\zeta_2'$ dan seterusnya,
tetapi ini hanya untuk berjaga-jaga jika Anda tidak mempercayai intuisi Anda;  fakta bahwa
$\|\phi^{-1}(u)-\phi^{-1}(v)-T^{-1}(u-v)\|$ kecil dibandingkan dengan
$\|u-v\|$ jelas merupakan akibat hipotesis bahwa
$\|\phi(x)-\phi(y)-T(x-y)\|$ kecil dibandingkan dengan $\|x-y\|$.

Argumen 263D sendiri sekarang tinggal memecah himpunan $D$ menjadi
bagian-bagian yang tepat, yang pada masing-masing bagiannya $\phi$ cukup
mendekati linear agar 263C dapat diterapkan, sehingga

\Centerline{$\mu^*\phi[D]\le\sum_{m=0}^{\infty}\mu^*\phi[D_m]
\le\sum_{m=0}^{\infty}(J_m+\epsilon)\mu^*D_m$.}

\noindent Dengan sedikit kehati-hatian (yang dilakukan dalam 263C, dengan syarat (i)),
kita juga dapat memastikan bahwa Jacobian $J$
dihampiri dengan baik oleh $J_m$ hampir di mana-mana dalam $D_m$, sehingga
$\sum_{m=0}^{\infty}J_m\mu^*D_m\bumpeq\int_DJ(x)dx$.

Gagasan-gagasan ini, jika digabungkan dengan hasil-hasil \S262, membawa kita
pada

\Centerline{$\int_EJ\,d\mu=\mu\phi[E]$}

\noindent ketika $\phi$ injektif dan $E\subseteq D$ terukur.
Kita memerlukan satu kiat terakhir, yang melibatkan himpunan Borel, untuk
mengubah ini menjadi

\Centerline{$\int_{\phi^{-1}[F]}J\,d\mu=\mu F$}

\noindent kapan pun $F\subseteq\phi[D]$ terukur, yang merupakan hal yang
diperlukan untuk penerapan 235J.

Saya berharap Anda sudah lama melihat, dan merasa senang oleh, cara dalam 263G.
Sekali lagi, ini sebenarnya bukan integrasi Lebesgue;  tetapi saya memasukkannya
untuk menunjukkan bahwa perangkat bab ini dapat digunakan menangani
hasil-hasil klasik, dan bahwa sesungguhnya kita memperoleh sedikit keuntungan
dari jerih payah kita, karena tidak perlu meminta maaf atas batas himpunan
$D$ yang menjadi sasaran sistem koordinat polar dalam pemetaan bidang.
Hasil yang sebenarnya telah saya berikan sebagai latihan dalam 252Xi.
Latihan itu melibatkan (jika Anda mengikuti semua acuannya) substitusi
berdimensi satu (yang dilakukan dalam 225Xh), teorema Fubini, dan penerapan
rumus-rumus \S235;  yakni, unsur-unsur yang hampir sama dengan yang digunakan
di atas, meskipun dalam urutan yang berbeda.   Saya dapat menyajikan ini tanpa
menyinggung diferensiasi dalam dimensi yang lebih tinggi karena perubahan
variabel pertama berdimensi satu, dan yang kedua (yang melibatkan fungsi
$x\mapsto\|x\|$, dalam 252Xi(i)) merupakan jenis yang sangat sederhana, sehingga
metode lain dapat digunakan untuk menemukan fungsi $J$.

Fungsi $R(y)=\sum_{x\in\phi^{-1}[\{y\}]}\sgn\det T(x)$ dalam 263Ib
termasuk dalam gagasan geometri diferensial yang lebih mendalam daripada
yang ingin saya uraikan di sini.   Dalam konteks berdimensi satu, fungsi itu
hanya menghitung perlintasan naik dan turun pada $y$ (lihat bagian (c)
pembuktian 263J), karena kita dapat memandang setiap $T(x)$ sebagai skalar
yang positif atau negatif.   (Relevan di sini bahwa dalam kasus ini kita dapat
mengandaikan bahwa $T(x)$ nonsingular kapan pun $\phi(x)=y$.)
Dalam dimensi yang lebih tinggi, saya kira hal pertama yang perlu dicari adalah
tafsiran geometris bagi $\sgn\det T(x)$.   (Lihat 263Ye.)
Tafsiran geometris untuk jumlahnya merupakan perkara lain lagi.
Namun patut diperhatikan bahwa (dengan tafsiran yang wajar bagi
"0$\times$tidak terdefinisi") kita dapat
menghubungkan $\int_Dg\phi\times\det\phi'$ dengan $\int_{\phi[D]}g\times R$,
dengan $R$ dapat didefinisikan dari $\phi$ tanpa mengacu pada $g$;  fungsi itu
menghitung lipatan dalam grafik $\phi$.

Gagasan abstrak yang menjadi pokok risalah ini memang
tidak membawa kita kepada banyak contoh khusus untuk melatih gagasan-gagasan
dalam bagian ini.   Contoh-contoh yang muncul cenderung sangat langsung,
seperti dalam 263G, 263Xa-263Xb, dan 263Xe.   Saya menyebut yang terakhir
karena contoh itu memberikan rumus yang diperlukan untuk membahas jenis baru
konvolusi (263Ya).   Pada hakikatnya, hal ini bergantung pada grup perkalian
$\Bbb R\setminus\{0\}$ sebagai pengganti grup penjumlahan $\Bbb R$ yang dibahas
dalam \S255.   Rumus $\bover1x$ dalam definisi $\nu$ tentu saja merupakan
turunan $\ln x$, dan $\ln$ merupakan isomorfisme antara
$(\ooint{0,\infty},\cdot,\nu)$ dan $(\Bbb R,+,\text{ukuran Lebesgue})$.
}%end of notes

\discrpage
